% 核对底本 PDF 第11页：知识体系·总论。人工逐节点转写，连线重新绘制。
\begin{center}
\begin{tikzpicture}[x=1cm,y=1cm,every node/.style={font=\fontsize{9.3}{12}\selectfont,inner sep=2pt,align=left},edge/.style={draw=black!65,line width=.45pt}]
\node[align=center] (root) at (0,5.4) {总\\论};
\draw[edge] (root.east)--(.5,5.4);\draw[edge] (.5,1.3)--(.5,10.5);
\draw[edge] (.5,10.5)--(.83,10.5);
\node[anchor=west,text width=14.2cm] at (.83,10.5) {刑法论：刑法的解释，刑法的原则，刑法的效力};
\node[anchor=west] (crime) at (.83,6.3) {犯罪论};\draw[edge] (.5,6.3)--(crime.west);\draw[edge] (crime.east)--(2.35,6.3);
\draw[edge] (2.35,2.8)--(2.35,9.6);
\foreach \yy in {9.6,8.75,7.9,7.05,6.2,5.35,4.5,3.65,2.8}{\draw[edge] (2.35,\yy)--(2.65,\yy);}
\node[anchor=west,text width=12.9cm] at (2.65,9.6) {\textbf{1. 犯罪客体：}法益};
\node[anchor=west,text width=12.9cm] at (2.65,8.75) {\textbf{2. 犯罪主体：}自然人；单位犯罪};
\node[anchor=west,text width=12.9cm] at (2.65,7.9) {\textbf{3. 客观要件：}危害行为；危害结果；因果关系};
\node[anchor=west,text width=12.9cm] at (2.65,7.05) {\textbf{4. 主观要件：}犯罪故意，犯罪过失；事实认识错误};
\node[anchor=west,text width=12.9cm] at (2.65,6.2) {\textbf{5. 排除违法事由：}正当防卫，紧急避险，被害人承诺};
\node[anchor=west,text width=12.9cm] at (2.65,5.35) {\textbf{6. 排除责任事由：}责任年龄，责任能力，违法性认识可能性，期待可能性};
\node[anchor=west,text width=12.9cm] at (2.65,4.5) {\textbf{7. 犯罪形态：}犯罪预备，犯罪未遂，犯罪中止，犯罪既遂};
\node[anchor=west,text width=12.9cm] at (2.65,3.65) {\textbf{8. 共同犯罪：}共同正犯，间接正犯，教唆犯，帮助犯，共同犯罪的特殊问题};
\node[anchor=west,text width=12.9cm] at (2.65,2.8) {\textbf{9. 罪数：}一个行为，两个行为};
\node[anchor=west] (penalty) at (.83,1.3) {刑罚论};\draw[edge] (.5,1.3)--(penalty.west);\draw[edge] (penalty.east)--(2.35,1.3);
\draw[edge] (2.35,.45)--(2.35,2.05);
\foreach \yy in {2.05,1.25,.45}{\draw[edge] (2.35,\yy)--(2.65,\yy);}
\node[anchor=west,text width=12.9cm] at (2.65,2.05) {\textbf{1. 刑罚的体系：}主刑，附加刑，非刑罚处罚措施};
\node[anchor=west,text width=12.9cm] at (2.65,1.25) {\textbf{2. 刑罚的裁量：}量刑情节，累犯，自首，立功，数罪并罚，缓刑};
\node[anchor=west,text width=12.9cm] at (2.65,.45) {\textbf{3. 刑罚的执行和消灭：}减刑，假释，时效};
\end{tikzpicture}
\end{center}
